\section{Vérification, validation et quantification de l'incertitude}
La théorie entourant le dénombrement ou l'estimation de quantités d'intérêts dans des populations finies est une sous-branche de la statistique. Initialement utilisées par les États pour dénombrer la population ou la production de denrées pendant l'Antiquité et le Moyen-Âge \parencite{oriol_formation_2007,tille_theorie_2019,clement_introduction_2020}. La pratique moderne trouve son origine dans les écrits de Laplace \parencite{laplace_theorie_1820}. La statistique a la particularité d'être imbriqué dans l'ensemble des autres disciplines scientifiques des sciences de la vie au génie qui est maintenant institutionnalisé dans la pratique et les institutions \parencite{national_research_council_assessing_2012}.\par
\textcite{national_research_council_assessing_2012} et \textcite{national_research_council_behavioral_2008} donnent un aperçu des méthodes couramment utilisées pour vérifier, valider et quantifier l'incertitude de modèles utilisées dans les sciences de la vie et les sciences sociales respectivement. Les définitions suivantes proviennent du deuxième rapport.
\begin{definition}
    La vérification est le processus pour déterminer si l'implémentation d'un modèle représente la description conceptuelle et les spécifications du concepteur.
\end{definition}
\begin{definition}
    La validation est le processus de détermination de la fidélité du modèle au monde réel de la perspective des usagers du modèle.
\end{definition}
Typiquement, la vérification serait assurée à l'aide de test unitaires sur les différentes fonctions pour assurer que le code fonctionne tel qu prévu. Dans le cadre de ce mémoire, la vérification a été faite de manière ad hoc. Les lecteurs peuvent se référer aux ouvrages de référence pertinents \parencite{sale_testing_2014,saleh_javascript_2013}.\par
En ce qui concerne la validation, le processus de validation consiste en diverses expériences ou observations pour déterminer à quel point le modèle est apte à prédire la quantité d'intérêt. Typiquement, ce processus est fait à différent niveaux hiérarchiques dans la simulation, idéalement au niveau le plus fondamental possible pour déterminer la précision du modèle au niveau systémique. Un processus de calibration ou des variables intermédiaires inconnues (que ce soit dans la littérature ou d'expériences précédentes) peut être utilisé pour caler une partie du modèle sur les résultats expérimentaux. \par
Le processus de validation peut être miné par l’inexactitude des mesures physiques utilisées pour comparer le modèle à la réalité, compliquant la tâche.